\chapter{Conclusion}
\label{chap:conclusion_and_future_work}
The trend since the early 2000s has been to move towards negative polarity, there was no explanation
whether the trend was due to efficiency or design factors. This thesis aimed on clarifying the reasons 
on why this was happening. After investigating 34 different subjects in common scenarios such as wiki-based
 and e-commerce the data revealed no significant difference between the two polarities. Therefore the conclusion 
 of this thesis is that there are not enough evidence to support the advantages of positive polarity over negative polarity.
However it is worth mentioning that users do prefer positive polarity over negative polarity, accepting the hypothesis H3 of this project.

\section{Future Work}
The future work of this project shoudl be focused in the following ideas:
\begin{itemize}
    \item \textbf{Sample Size:} The sample size of this project is short, increasing the sample may reveal significant results.
    \item \textbf{Precise Regression Study:} The regressions were calculated with a rough calculation, there may be tools
    that provide this variable with more precision.
    \item \textbf{Longer Exposure:} The experiment is centered around texts however, the expsure and time spent reading may reveal new 
    results by expanding the time of the experiment.
    \item \textbf{Sample Age:} The sample of the project is focused on young people, they have been adapted and exposed to new technologies
    since they were born, providing a older sample may reveal new results as older people have not adapted entirely to the internet and its design.
\end{itemize}
